% Separate analytical review increment. No earlier proof file is modified.
% Unique prefix: inc:AN02:tail:v1:
\documentclass[11pt]{article}
\usepackage[T1]{fontenc}
\usepackage[utf8]{inputenc}
\usepackage{lmodern}
\usepackage[margin=0.9in]{geometry}
\usepackage{amsmath,amssymb,amsthm,mathtools,booktabs,array,microtype}
\usepackage{xurl}
\usepackage[hidelinks]{hyperref}
\usepackage{fancyhdr}
\hypersetup{pdftitle={AN02: Strong ancestry and nonlinear upper-tail transport},pdfauthor={Research proof increment for independent review}}
\pagestyle{fancy}\fancyhf{}
\fancyhead[L]{\small AN02/AN03: ancestry and upper-tail transport}
\fancyhead[R]{\small Review increment v1}
\fancyfoot[C]{\thepage}
\setlength{\headheight}{14pt}
\setlength{\emergencystretch}{3em}
\allowdisplaybreaks[2]
\numberwithin{equation}{section}
\newtheorem{theorem}{Theorem}[section]
\newtheorem{lemma}[theorem]{Lemma}
\newtheorem{proposition}[theorem]{Proposition}
\newtheorem{corollary}[theorem]{Corollary}
\theoremstyle{remark}\newtheorem{remark}[theorem]{Remark}
\newcommand{\E}{\mathbb E}
\newcommand{\R}{\mathbb R}
\newcommand{\dd}{\,d}
\newcommand{\1}{\mathbf 1}
\newcommand{\scope}[1]{\par\noindent\textbf{Scope.} #1\par\smallskip}
\newcommand{\gap}[1]{\par\noindent\textbf{Not established.} #1\par\smallskip}
\title{Strong-label ancestry and nonlinear upper-tail transport\\
\large An initialized mass bound, an exact outer comparison, and resident-passage interfaces}
\author{Research proof increment for independent review}
\date{Version 1 --- October 2, 2026}
\begin{document}
\maketitle
\begin{abstract}
We prove an original finite-noise, initialized measure lower bound that retains
both positive initial-coordinate weights. In particular, every fixed ancestral
subset of the initial right half-circle retains its $\cos^2\alpha$-weighted
strong amplification through the established early comparison interval. We
then solve the zero-noise target-only comparison exactly, including its
finite-time radial correction, and compute the one-sided cubic ancestry
tail and its angular pushforward. This comparison predicts a tail exponent
$3(1-\rho^2)/2=5/6$, not $4/3$, at $\rho=2/3$; it is not a captured-seed
law for the original flow. Within the leading scaled system, we prove an
exact nonlinear upper-diagonal test-particle equation along the resident
strong-learning curve. Every positive diagonal displacement grows after
saturation; the far-field linear intercept $a/\rho^2$ is not a separatrix.
For a general self-consistent leading population, a separate same-field
comparison propagates an entire upper-tail profile and bounds the mass that
exits a prescribed core neighborhood during a resident interval. Finally,
we state source-resolved reaction/flux identities and the half-mass-centered
AN3/AN7 contract. The coupled initialized strong-learning and weak-capture
lemma, and therefore initialized Theorems A and B, remain unproved.
\end{abstract}

\section{Contract, dependencies, and distinctions}
\label{inc:AN02:tail:v1:contract}
The frozen base is \path{THEOREM_A_ANALYTICAL_PROGRESS.tex}, version 1.0,
and \path{RELU_POPULATION_FOUNDATIONS_v2.tex}. Dependencies are
\texttt{pa:targetcomparison} and \texttt{pa:earlybounds} for Section 2;
the Gaussian feature identity in foundations P4 for the target-only field;
and \texttt{pa:leading}, \texttt{pa:masslaws}, \texttt{pa:resident}, and
\texttt{pa:overlap} for Sections 4--6. We restate the equations used.
The witness target remains the reviewed
\path{AN06_learned_orbit_witnesses_v2.tex}; it is not revised here.

Three evolutions must remain distinct:
\begin{enumerate}
\item The original $q>0$, self-consistent, isotropically initialized population.
\item An explicitly defined \emph{zero-noise target-only comparison}, which
has neither Gaussian overlap nor the learned residual feedback.
\item The defined leading scaled population and, only in Section 4, passive
characteristics in its prescribed coherent resident field.
\end{enumerate}
No statement transfers (2) or a prescribed-field result in (3) to (1)
without additional estimates. In particular, initial-label ancestry is not
present-time family membership, a tail at saturation is not its mass after
resident passage, and a passive tail is not a finite-mass self-consistent
bridge. The new same-field result in Section 5 avoids the last substitution,
but requires its own entry and reference-characteristic hypotheses.

Use $\lambda=\rho^2\in(0,1)$, normalized time $\tau$, and
$\lambda_0(\dd\alpha)=\dd\alpha/(2\pi)$. Put
\[
 P_1=\mathcal N(e_1,q^2 I),\quad P_2=\mathcal N(\rho e_2,q^2I),\quad
 P=(P_1+P_2)/2,\qquad
 K(r)=\phi(r)+r\Phi(r).
\]
The original initialization is $U_0=W_0=\sqrt{S_0}a_\alpha$ and $m_0=S_0$.
For a fixed label set $B$, write
\[
 \nu_\tau^{\rm lab}(B)=\int_Bm_\tau(\alpha)\dd\lambda_0(\alpha).
\]
This is a measure on \emph{initial labels}, not the angular pushforward.
All proofs below are analytical. Compilation is a typesetting check, not
a verification of the mathematical conclusions.

\section{An initialized positive-coordinate measure lower bound}
\label{inc:AN02:tail:v1:ancestry}
Let $\widehat U$ denote the exact finite-$q$ target-only comparison:
\begin{equation}\label{inc:AN02:tail:v1:target}
 \widehat U'=\E_P[X(\widehat U^\top X)_+],\qquad
 \widehat U(0,\alpha)=\sqrt{S_0}a_\alpha,\qquad
 \widehat m=|\widehat U|^2.
\end{equation}
Set $a_*=1/2+q^2$, and define the already-proved comparison quantities
\begin{equation}\label{inc:AN02:tail:v1:comparison}
 \overline S=S_0e^{2a_*\tau},\quad
 \epsilon_{\rm al}=\overline S-S_0,\quad
 H=S_0(2e^{2a_*\tau}-3e^{a_*\tau}+1).
\end{equation}
On every interval with $\epsilon_{\rm al}<\pi$, the checkpoint gives
\begin{equation}\label{inc:AN02:tail:v1:masscomparison}
 S\le\overline S,\quad
 |\theta-\widehat\theta|\le H,\quad
 |\psi-\widehat\theta|\le H+\epsilon_{\rm al},\quad
 \left|\log\frac{m}{\widehat m}\right|\le2H.
\end{equation}

\begin{theorem}[Positive-coordinate ancestral measure domination]
\label{inc:AN02:tail:v1:weightedmass}
For every $q>0$, $0<\rho<1$, $S_0>0$, and time with
$\epsilon_{\rm al}(\tau)<\pi$, the original initialized population satisfies
\begin{equation}\label{inc:AN02:tail:v1:pointweight}
 m_\tau(\alpha)\ge e^{-2H(\tau)}S_0
 \left[e^\tau(\cos\alpha)_+^2+e^{\lambda\tau}(\sin\alpha)_+^2\right].
\end{equation}
Consequently, for every measurable fixed initial-label set $B$,
\begin{equation}\label{inc:AN02:tail:v1:measureweight}
 \nu_\tau^{\rm lab}(B)\ge e^{-2H(\tau)}S_0
 \int_B\left[e^\tau(\cos\alpha)_+^2+
 e^{\lambda\tau}(\sin\alpha)_+^2\right]\dd\lambda_0.
\end{equation}
In particular,
\begin{equation}\label{inc:AN02:tail:v1:totalweight}
 S(\tau)\ge\frac{e^{-2H(\tau)}S_0}{4}(e^\tau+e^{\lambda\tau}),
 \qquad
 \nu_\tau^{\rm lab}\{\cos\alpha>0\}\ge
 \frac{e^{-2H(\tau)}S_0e^\tau}{4}.
\end{equation}
\end{theorem}
\scope{Original finite-noise flow, from the prescribed initialization, on
its proved early comparison interval. No captured-family interpretation is assumed.}
\begin{proof}
Write $r=|\widehat U|>0$, $\widehat a=\widehat U/r$, and
$\mu_1=1$, $\mu_2=\rho$. Gaussian integration gives, for $j=1,2$,
\begin{equation}\label{inc:AN02:tail:v1:coordinatefield}
 \widehat U_j'=\frac{\mu_jqr}{2}
 K\left(\frac{\mu_j\widehat a_j}{q}\right)
 +d_q(\widehat a)\widehat U_j,\qquad
 d_q(\widehat a)=\frac{q^2}{2}
 \left[\Phi\left(\frac{\widehat a_1}{q}\right)
 +\Phi\left(\frac{\rho\widehat a_2}{q}\right)\right]\ge0.
\end{equation}
For completeness, for a Gaussian with mean $\mu_j e_j$, Stein's identity
reads
\[
 \E[X(\widehat U^\top X)_+]
 =\mu_j e_j\,qrK(\mu_j\widehat a_j/q)
 +q^2\widehat U\Phi(\mu_j\widehat a_j/q).
\]
Averaging the two clusters proves \eqref{inc:AN02:tail:v1:coordinatefield}.
The covariance term from \emph{both} clusters is retained.

An initially positive coordinate stays positive: at a possible zero its
vector field is $\mu_jqr\phi(0)/2>0$. The radius cannot vanish; for the
target-only flow $(|\widehat U|^2)'=2\E(\widehat U^\top X)_+^2\ge0$.
On a positive coordinate, $K(z)\ge z$ and $d_q\ge0$ give
\[
 \widehat U_j'\ge\frac{\mu_j^2}{2}\widehat U_j,
 \qquad
 \widehat U_j(\tau)\ge\widehat U_j(0)e^{\mu_j^2\tau/2}.
\]
Squaring and summing only the coordinates that were positive initially
proves the bound for $\widehat m$. Multiplying by $e^{-2H}$, as allowed
by \eqref{inc:AN02:tail:v1:masscomparison}, proves
\eqref{inc:AN02:tail:v1:pointweight}. Integration is over fixed labels.
Finally $\int(\cos\alpha)_+^2\dd\lambda_0=
\int(\sin\alpha)_+^2\dd\lambda_0=1/4$.
\end{proof}

\begin{corollary}[A nonvanishing early strong-ancestry mass fraction]
\label{inc:AN02:tail:v1:earlyfraction}
Let $0<S_0<\delta\le1/100$ and
\[
 T_\delta=\frac{\log(\delta/S_0)}{1+2q^2}.
\]
Then $S\le\delta$ through $T_\delta$ and
\begin{equation}\label{inc:AN02:tail:v1:earlyfractioneq}
 \nu_{T_\delta}^{\rm lab}\{\cos\alpha>0\}
 \ge\frac{e^{-4\delta}\delta}{4}
 \exp\left[-\frac{2q^2}{1+2q^2}\log\frac\delta{S_0}\right].
\end{equation}
If $q^2\log(\delta/S_0)\le A$, the right side is at least
$\delta e^{-4\delta-2A}/4$.
\end{corollary}
\begin{proof}
The endpoint comparison is
$H(T_\delta)=2\delta-3\sqrt{S_0\delta}+S_0\le2\delta$.
Also $S_0e^{T_\delta}=\delta e^{-2q^2T_\delta}$.
Apply \eqref{inc:AN02:tail:v1:totalweight}.
\end{proof}
\gap{The mass in \eqref{inc:AN02:tail:v1:earlyfractioneq} has not been shown
to occupy a strong scaled chart or to fit the strong concept. This improves
the initialized mass input, not the missing order-one learning continuation.}

\subsection{The inherited tail weight is explicit}
For $R>0$ define a fixed ancestral subset
\[
 B_R=\{\alpha\in(0,\pi/2):\tan\alpha\ge R\}.
\]
Put
\begin{align}
 I_c(R)&=\int_{\arctan R}^{\pi/2}\cos^2\alpha\dd\alpha
 =\frac12\left[\arctan(1/R)-\frac{R}{1+R^2}\right],\label{inc:AN02:tail:v1:Ic}\\
 I_s(R)&=\int_{\arctan R}^{\pi/2}\sin^2\alpha\dd\alpha
 =\frac12\left[\arctan(1/R)+\frac{R}{1+R^2}\right].\label{inc:AN02:tail:v1:Is}
\end{align}
Theorem~\ref{inc:AN02:tail:v1:weightedmass} yields the exact original-flow
lower bound
\begin{equation}\label{inc:AN02:tail:v1:ancestraltail}
 \nu_\tau^{\rm lab}(B_R)\ge
 \frac{e^{-2H(\tau)}S_0e^\tau}{2\pi}
 [I_c(R)+e^{-(1-\lambda)\tau}I_s(R)].
\end{equation}
Substitution $t=\tan\alpha$ gives
\[
 I_c(R)=\int_R^\infty\frac{\dd t}{(1+t^2)^2},\qquad
 I_s(R)=\int_R^\infty\frac{t^2\dd t}{(1+t^2)^2}.
\]
For $R\ge1$, pointwise comparison with powers of $t$ proves
\begin{equation}\label{inc:AN02:tail:v1:Itailbounds}
 \frac{1}{3R^3(1+R^{-2})^2}\le I_c(R)\le\frac{1}{3R^3},\qquad
 \frac{1}{R(1+R^{-2})^2}\le I_s(R)\le\frac1R.
\end{equation}
In particular $R^3I_c(R)\to1/3$ and $RI_s(R)\to1$ by these inequalities.
The finite-time $I_s$ term must not be silently deleted. This is an
ancestral measure lower bound; no current-angle lower bound follows from it.

\section{Exact outer transport in the zero-noise target-only comparison}
\label{inc:AN02:tail:v1:outer}
\scope{An auxiliary comparison, not the original positive-noise or learned
population. Here zero noise is built into the definition, not taken through
an unproved long-time limit.}
Define $Z$ by
\begin{equation}\label{inc:AN02:tail:v1:zerofield}
 Z'=\frac12(Z_1)_+e_1+\frac\lambda2(Z_2)_+e_2,
 \qquad Z(0,\alpha)=\sqrt{S_0}a_\alpha.
\end{equation}
This is the exact target-only field for the two point masses
$e_1,\rho e_2$ with equal weights. It is globally Lipschitz, including
coordinate axes. Write $\theta^{(0)}=\arg Z$ and $m^{(0)}=|Z|^2$.

\begin{proposition}[Quadrant solution and finite-time mass weighting]
\label{inc:AN02:tail:v1:quadrants}
Let $r_+(\tau)=e^{-(1-\lambda)\tau/2}$ and
$r_-(\tau)=e^{-\tau/2}$. In the open first quadrant,
\begin{equation}\label{inc:AN02:tail:v1:Q1}
 \tan\theta^{(0)}=r_+\tan\alpha,\qquad
 m^{(0)}=S_0e^\tau(\cos^2\alpha+r_+^2\sin^2\alpha).
\end{equation}
In the open fourth quadrant, with $-\pi/2<\alpha<0$,
\begin{equation}\label{inc:AN02:tail:v1:Q4}
 \tan\theta^{(0)}=r_-\tan\alpha,\qquad
 m^{(0)}=S_0e^\tau(\cos^2\alpha+r_-^2\sin^2\alpha).
\end{equation}
The initial right-half-circle mass is exactly
\begin{equation}\label{inc:AN02:tail:v1:rightmass}
 M_{\rm RH}^{(0)}(\tau)=\frac{S_0e^\tau}{8}(2+r_+^2+r_-^2).
\end{equation}
Its normalized initial-label measure converges, in density in $L^1$, to
$(2/\pi)\cos^2\alpha\dd\alpha$ on $(-\pi/2,\pi/2)$ as $\tau\to\infty$.
\end{proposition}
\begin{proof}
Every initially positive coordinate has its scalar exponential growth, and
every initially negative coordinate is constant. In Q1,
$Z=\sqrt{S_0}(e^{\tau/2}\cos\alpha,e^{\lambda\tau/2}\sin\alpha)$;
in Q4 its second coordinate is $\sqrt{S_0}\sin\alpha$.
This proves the two angle/mass formulas. Integrate using
$\int_0^{\pi/2}\cos^2=\int_0^{\pi/2}\sin^2=\pi/4$.
The normalized densities converge uniformly separately on each quadrant,
hence in $L^1$ on their union. Coordinate-axis labels have zero initial
Lebesgue measure; no boundary mass is omitted in the integrals.
\end{proof}

\begin{proposition}[One-sided angular tails at every finite time]
\label{inc:AN02:tail:v1:angulartail}
For $0<\vartheta<\pi/2$ and $b=\tan\vartheta$, the Q1 mass at
$\theta^{(0)}\ge\vartheta$ is
\begin{equation}\label{inc:AN02:tail:v1:tailpush}
 \mathcal T_+^{(0)}(\tau,\vartheta)
 =\frac{S_0e^\tau}{2\pi}
 \left[I_c(b/r_+)+r_+^2I_s(b/r_+)\right].
\end{equation}
The Q4 mass at $\theta^{(0)}\le-\vartheta$ is the same expression with
$r_-$ in place of $r_+$. Thus at fixed $b>0$, as $\tau\to\infty$,
\begin{equation}\label{inc:AN02:tail:v1:asymmetricratio}
 \frac{\mathcal T_+^{(0)}(\tau,\vartheta)}
 {\mathcal T_-^{(0)}(\tau,\vartheta)}
 \sim e^{3\lambda\tau/2}.
\end{equation}
\end{proposition}
\begin{proof}
In Q1 the terminal angular event is exactly
$\tan\alpha\ge b/r_+$. Substitute \eqref{inc:AN02:tail:v1:Q1} in its
mass integral. Q4 follows by $\alpha\mapsto-\alpha$.
At fixed $b$, \eqref{inc:AN02:tail:v1:Itailbounds} gives
\[
 I_c(b/r)+r^2I_s(b/r)\sim r^3\left(\frac1{3b^3}+\frac1b\right).
\]
Take the ratio and use $r_+/r_-=e^{\lambda\tau/2}$.
\end{proof}

\begin{corollary}[The reference tail exponent at a nominal strong scale]
\label{inc:AN02:tail:v1:exponent}
Fix $A>0$, $0<\rho<1$, and $\zeta>0$. Set
$\tau_A=\log(A/S_0)$, so $S_0e^{\tau_A}=A$; this is a
\emph{comparison} clock, not a true saturation time. As $q\downarrow0$,
$S_0\downarrow0$, suppose
\[
 \frac{(S_0/A)^{(1-\lambda)/2}}q\longrightarrow0.
\]
Then
\begin{equation}\label{inc:AN02:tail:v1:tailpower}
 \mathcal T_+^{(0)}(\tau_A,q\zeta)
 \sim\frac{A}{6\pi\zeta^3}
 \left(\frac{S_0}{A}\right)^{3(1-\lambda)/2}q^{-3},
\end{equation}
whereas its fraction of $M_{\rm RH}^{(0)}$ is asymptotic to
\begin{equation}\label{inc:AN02:tail:v1:tailfraction}
 \frac{2}{3\pi\zeta^3}
 \left(\frac{S_0}{A}\right)^{3(1-\lambda)/2}q^{-3}.
\end{equation}
At $\rho=2/3$, the initialization exponent is $5/6$.
In particular the mass in \eqref{inc:AN02:tail:v1:tailpower}, divided by
$S_0$, diverges in this regime. It is not $o(S_0)$.
\end{corollary}
\begin{proof}
Put $r=(S_0/A)^{(1-\lambda)/2}$ and $b=\tan(q\zeta)$.
Then $b\sim q\zeta$, $r/b\to0$, and
\[
 \frac{r^2I_s(b/r)}{I_c(b/r)}\sim3b^2\to0.
\]
Use \eqref{inc:AN02:tail:v1:tailpush} and
$I_c(b/r)\sim r^3/(3b^3)$. Equation
\eqref{inc:AN02:tail:v1:rightmass} is asymptotic to $A/4$.
Finally $3(1-4/9)/2=5/6$; the ratio to $S_0$ has the positive
constant factor times $S_0^{-1/6}q^{-3}$.
\end{proof}
\gap{The reference tail power is not a proof of either an actual tail mass
at strong saturation or of later weak recruitment. The exponent correction
invalidates a subdominance inference from that proposed power, not AN02's
proved reservoir lower bound. More recruited weak mass could help the total
seed lower bound, but its angular law and feedback must then be controlled.}

\paragraph{Regime clocks in this comparison.}
For a fixed Q1 label with $\tan\alpha>0$, the first time at
$\tan\theta^{(0)}=qR$ is exactly
\[
 \tau_c=\frac{2}{1-\lambda}\log\frac{\tan\alpha}{qR}
\]
when the logarithm is positive. For Q4 the coefficient is $2$.
These are exact individual comparison hitting times, not a uniform
charting theorem for the whole isotropic population. Labels arbitrarily
close to $\pi/2$ always require separate tail accounting.

\section{Nonlinear upper-tail motion in the resident field}
\label{inc:AN02:tail:v1:resident}
\scope{Exact passive-characteristic statements within the defined leading
model, in a prescribed coherent residual field. They do not assert that a
finite-mass tail leaves that residual field unchanged.}
Let $a=a_\rho>0$ solve $a=\rho K(\rho a)$ and put
\[
 P=\Phi(\rho a),\qquad \alpha_{\rm sh}=\frac\lambda2P.
\]
The notation $\alpha_{\rm sh}$ denotes the checkpoint's shape rate $\alpha$,
not an initial label. The prescribed resident solution is
$\mu_s=M(\tau)\delta_{(a,a)}$, $\mu_w=0$, $M'=M(1-M)$,
with $0<M\le1$. A passive strong characteristic obeys
\begin{align}
 2x'&=-(1-M)x+\rho\phi(\rho x)-\rho M F(\rho x,\rho a)
 +\lambda y\Phi(\rho x),\label{inc:AN02:tail:v1:passive}\\
 2y'&=-(1-M)y-Ma+\rho K(\rho x),\nonumber
\end{align}
where $F(t,r)=\phi(\min(t,r))+r\Phi(\min(t,r))$.

\begin{theorem}[Exact nonlinear diagonal equation]
\label{inc:AN02:tail:v1:nonlineardiagonal}
If $x(\tau_0)=y(\tau_0)=a+h_0$, $h_0>0$, then
$x=y=a+h$ with $h>0$, and
\begin{equation}\label{inc:AN02:tail:v1:heq}
 h'=\left[-\frac{1-M}{2}+\mathcal B(h)\right]h,\qquad
 \mathcal B(h)=\frac{\rho K(\rho(a+h))-a}{2h}
 =\frac\lambda{2h}\int_0^h\Phi(\rho(a+r))\dd r.
\end{equation}
The function $\mathcal B$ is strictly increasing, with
\begin{equation}\label{inc:AN02:tail:v1:Bhbounds}
 \alpha_{\rm sh}<\mathcal B(h)<\frac\lambda2\quad(h>0),\qquad
 \lim_{h\downarrow0}\mathcal B(h)=\alpha_{\rm sh},\quad
 \lim_{h\to\infty}\mathcal B(h)=\frac\lambda2.
\end{equation}
For every finite $\tau\ge\tau_0$,
\begin{equation}\label{inc:AN02:tail:v1:nonlinearmultipliers}
 h_0e^{\alpha_{\rm sh}(\tau-\tau_0)}
 \sqrt{\frac{M(\tau_0)}{M(\tau)}}
 \le h(\tau)\le
 h_0e^{\lambda(\tau-\tau_0)/2}
 \sqrt{\frac{M(\tau_0)}{M(\tau)}}.
\end{equation}
At a state of displacement $h$, growth begins when
$M>1-2\mathcal B(h)$, a threshold strictly between
$1-\lambda$ and $1-2\alpha_{\rm sh}$.
\end{theorem}
\begin{proof}
For $x\ge a$, $F(\rho x,\rho a)=K(\rho a)=a/\rho$.
On $x=y=z\ge a$, both equations in
\eqref{inc:AN02:tail:v1:passive} therefore equal
$[-(1-M)z-Ma+\rho K(\rho z)]/2$.
Uniqueness makes this half-diagonal invariant. Subtract $a$ and use the
resident equation. The integral expression for $\mathcal B$ follows from
$[\rho K(\rho z)]'=\lambda\Phi(\rho z)$. An average of a strictly
increasing function over $[0,h]$ strictly increases with $h$, proving the
monotonicity and bounds. Its limit at infinity is $\lambda/2$ by the
limit of $\Phi$; its limit at zero follows by continuity.
The equation is locally Lipschitz at $h=0$, and the bounds keep $h$
positive and prevent finite-time blowup. Integrating
$\alpha_{\rm sh}-(1-M)/2\le(\log h)'\le\lambda/2-(1-M)/2$
and $\int(1-M)=\log[M(\tau)/M(\tau_0)]$ proves the multipliers.
The growth threshold is the sign condition in \eqref{inc:AN02:tail:v1:heq}.
\end{proof}

\begin{corollary}[No second diagonal separatrix after saturation]
\label{inc:AN02:tail:v1:noseparatrix}
For $M\equiv1$, every $h_0>0$ grows and reaches every larger finite
$H>h_0$. Its exact hitting time is
\begin{equation}\label{inc:AN02:tail:v1:hitting}
 T(h_0,H)=2\int_{a+h_0}^{a+H}
 \frac{\dd z}{\rho K(\rho z)-a},\qquad
 \frac2\lambda\log\frac H{h_0}
 \le T(h_0,H)\le\frac1{\alpha_{\rm sh}}\log\frac H{h_0}.
\end{equation}
In particular $x=y=a/\lambda$ is not an equilibrium or separatrix:
\begin{equation}\label{inc:AN02:tail:v1:intercept}
 x'=y'=\frac\rho2\left[K(a/\rho)-a/\rho\right]>0
 \quad\text{there}.
\end{equation}
\end{corollary}
\begin{proof}
Use \eqref{inc:AN02:tail:v1:heq} with $M=1$ and integrate the scalar
positive velocity. For the last formula, $\rho x=a/\rho$ and
$\rho K(\rho x)-a=\rho[K(a/\rho)-a/\rho]>0$.
\end{proof}

\begin{remark}[The far-field linear intercept]
For $x\ge a$ and $M=1$, the exact field is
\[
 x'=\tfrac12[\rho\phi(\rho x)-a+\lambda y\Phi(\rho x)],\qquad
 y'=\tfrac12[\rho K(\rho x)-a].
\]
Its large-$x$ linear approximation is
$x'\approx(\lambda y-a)/2$, $y'\approx(\lambda x-a)/2$.
The linear system's saddle at $(a/\lambda,a/\lambda)$ is an intercept of
that approximation, not an exact threshold of the field at finite $x$.
The exact local instability at $(a,a)$ and the far-field rate $\lambda/2$
are linked by \eqref{inc:AN02:tail:v1:heq}. No global two-dimensional
separatrix description for arbitrary off-diagonal particles is claimed.
\end{remark}

\paragraph{Passive tail crossing over an entire resident interval.}
With $M=1$ and a measure of passive diagonal displacements $h_0>0$,
let $h_{\rm crit}(T)$ be the unique displacement that reaches $H$ at time
$T>0$. Scalar order and \eqref{inc:AN02:tail:v1:hitting} give
\[
 He^{-\lambda T/2}\le h_{\rm crit}(T)\le He^{-\alpha_{\rm sh}T}.
\]
The mass that has crossed $H$ by $T$ is exactly the initial passive mass
on $\{h_0\ge h_{\rm crit}(T)\}$. It is bounded below by the mass on
$\{h_0\ge He^{-\alpha_{\rm sh}T}\}$ and above by the mass on
$\{h_0\ge He^{-\lambda T/2}\}$. These are passive-field bounds only.

\section{A self-consistent same-field upper-tail propagation theorem}
\label{inc:AN02:tail:v1:fulltail}
We now allow an arbitrary leading population. On a finite interval
$[\tau_0,T]$, assume a characteristic solution with bounded support,
$0<M\le1$, $0\le N\le1$, and
\[
 M'=M(1-M),\qquad N'=\lambda N(1-N).
\]
Use total moments $Y=\int y\dd\mu_s$, $V=\int v\dd\mu_w$,
$K_w=\int K(u)\dd\mu_w$, and
$\mathcal S(x)=\int F(\rho x,\rho\widetilde x)\dd\mu_s$.
The strong field is
\begin{align}
 2x'&=-(1-M)x+\rho\phi(\rho x)-\rho\mathcal S(x)
 +\lambda[V+(1-N)y]\Phi(\rho x),\label{inc:AN02:tail:v1:fullstrong}\\
 2y'&=-(1-M)y-Y-K_w+\rho(1-N)K(\rho x).\nonumber
\end{align}
Choose a reference characteristic $(x_c,y_c)$ solving this \emph{same}
field. It need not carry positive mass. It cannot be replaced by a
coherent orbit from a different field without an error estimate.
Let $\mathfrak m_\tau(\dd\beta)$ be the strong mass measure on its fixed
characteristic labels, so that $\mu_s(\tau)$ is its pushforward by
$\beta\mapsto(x(\tau,\beta),y(\tau,\beta))$. This does not require
that the coordinate map be injective.
Let $B$ be a fixed set of strong labels initially co-ordered above it:
$h_0=x(\tau_0)-x_c(\tau_0)\ge0$ and
$k_0=y(\tau_0)-y_c(\tau_0)\ge0$ on $B$.

Define the population tail moment, core imbalance, and its positive budget
\begin{equation}\label{inc:AN02:tail:v1:tailbudget}
 \Pi(\tau,r)=\int(\widetilde x-r)_+\dd\mu_s,\qquad
 U_c=V+(1-N)y_c-x_c,\qquad
 b_c=\frac{\rho^3\phi(0)}2(U_c)_+.
\end{equation}
For $H>0$ put
\begin{align}
 c_H(\tau)&=\frac\lambda2(1-N)\Phi(\rho(x_c+H)),&
 c_\infty(\tau)&=\frac\lambda2(1-N),\label{inc:AN02:tail:v1:cH}\\
 G_H(t)&=\exp\int_{\tau_0}^t
 \left[-\frac{1-M}2+c_H+b_c\right]\dd\tau,&
 G_H^*(T)&=\max_{\tau_0\le t\le T}G_H(t).\nonumber
\end{align}
Define $G_\infty$ by replacing $c_H$ with $c_\infty$.
For $r>0$ let
\[
 \mathcal P_0(r)=\int_B\1_{\{h_0+k_0\ge r\}}\dd\mathfrak m_{\tau_0}.
\]
This is an unnormalized upper-tail profile in initial strong mass.

\begin{theorem}[Nonlinear profile and first-exit mass bounds]
\label{inc:AN02:tail:v1:profiletheorem}
The labels in $B$ remain co-ordered above $(x_c,y_c)$. Put $D=h+k$.
Then
\begin{equation}\label{inc:AN02:tail:v1:globalD}
 D(t)\le G_\infty(t)D(\tau_0),\qquad
 \mathfrak m_t\{\beta\in B:D(t,\beta)\ge r\}
 \le\frac{M(t)}{M(\tau_0)}\mathcal P_0(r/G_\infty(t)).
\end{equation}
Let $E_H(T)$ be the initial labels in $B$ for which
$\sup_{\tau_0\le t\le T}D(t)\ge H$, including labels already outside
at $\tau_0$. Then
\begin{equation}\label{inc:AN02:tail:v1:exitmass}
 \mathfrak m_T(E_H(T))
 \le\frac{M(T)}{M(\tau_0)}
 \mathcal P_0\left(\frac H{G_H^*(T)}\right).
\end{equation}
If $\mathcal P_0(r)\le C(\varepsilon/r)^p$ for all $r>0$, then
\begin{equation}\label{inc:AN02:tail:v1:powerexit}
 \mathfrak m_T(E_H(T))\le\frac{M(T)}{M(\tau_0)}
 C\left(\frac{\varepsilon G_H^*(T)}H\right)^p.
\end{equation}
These are full same-field leading-population bounds, not passive-field
approximations. They do not themselves establish the profile hypothesis or
bound the output of labels after chart exit in the original Gaussian flow.
\end{theorem}
\begin{proof}
The strong cross derivatives are
$\partial_yx'=\partial_xy'=\lambda(1-N)\Phi(\rho x)/2\ge0$.
Local Lipschitz regularity follows from the stated kernel and bounded
support. The standard first-exit comparison, equivalently its strict
positive-barrier version followed by a limit, preserves $h,k\ge0$ in a
common time-dependent field. No sign on donor couplings between different
populations is used.

Differentiate the overlap with respect to its receiver coordinate:
\[
 \mathcal S'(r)=\rho^2\phi(\rho r)\Pi(\tau,r).
\]
This follows directly from $F_t(t,z)=(z-t)_+\phi(t)$; bounded support
justifies differentiation under the integral. Subtract the two
characteristic equations, evaluating the $y$-difference in the first
one at the upper receiver input. The \emph{exact} pair identities are
\begin{align}
 2h'={}&-(1-M)h+
 \rho^3\int_{x_c}^{x_c+h}
 [U_c-(r-x_c)-\Pi(\tau,r)]\phi(\rho r)\dd r\nonumber\\
 &+\lambda(1-N)k\Phi(\rho(x_c+h)),\label{inc:AN02:tail:v1:exactpair}\\
 2k'={}&-(1-M)k+
 \lambda(1-N)\int_{x_c}^{x_c+h}\Phi(\rho r)\dd r.\nonumber
\end{align}
Since $h\ge0$, both $r-x_c$ and $\Pi$ in the first integral are
nonnegative. Its contribution to $h'$ is at most $b_ch$.
With $\Phi\le1$, addition proves
\[
 D'\le[-(1-M)/2+c_\infty+b_c]D.
\]
When $D\le H$, one also has $h\le H$, so monotonicity of $\Phi$ gives
\[
 D'\le[-(1-M)/2+c_H+b_c]D.
\]
These inequalities do not divide by $h$, $k$, or $D$; zero differences
remain exactly zero. Gronwall gives the global bound, and the local bound
up to the first hit of $H$.

All fixed strong labels have the same relative reaction rate. Thus
$m(t,\beta)/m(\tau_0,\beta)=M(t)/M(\tau_0)$.
The event $D(t)\ge r$ implies $D(\tau_0)\ge r/G_\infty(t)$,
which proves the instantaneous tail estimate. A label that first hits $H$
at time $t_e$ satisfies
$H\le G_H(t_e)D(\tau_0)\le G_H^*(T)D(\tau_0)$.
Labels already outside at entry satisfy the same implication because
$G_H^*(T)\ge G_H(\tau_0)=1$. Integrate this event inclusion over initial
labels and multiply by $M(T)/M(\tau_0)$.
The power-profile consequence is direct substitution.
\end{proof}

\paragraph{Relation to the competing passage rates.}
For $N(\tau_0)>0$ the global amplification is exactly
\begin{equation}\label{inc:AN02:tail:v1:globalG}
 G_\infty(t)=
 \sqrt{\frac{M(\tau_0)}{M(t)}}
 \sqrt{\frac{N(t)}{N(\tau_0)}}
 \exp\int_{\tau_0}^t b_c\dd\tau.
\end{equation}
For $N\equiv0$, replace the second square root by
$e^{\lambda(t-\tau_0)/2}$. This follows by integrating the two logistic
laws. At weak half-mass the global tail estimate therefore pays
$(2N_{\rm ent})^{-1/2}$, not the resident-core factor
$(2N_{\rm ent})^{-\alpha_{\rm sh}/\lambda}$.
For a prescribed resident core $x_c=y_c=a$, $N=0$, one has $b_c=0$,
and the local coefficient $c_H=\lambda\Phi(\rho(a+H))/2$ tends to
$\alpha_{\rm sh}$ as $H\downarrow0$. This explains why a useful proof
must distinguish small core displacements from labels that leave the core;
a single global Lipschitz rate loses the desired exponent balance.
The exact pair identity retains a favorable upper-tail donor term through
$-\Pi$; the upper bounds above discard it conservatively.

\section{Source-resolved capture and the clock-to-witness contract}
\label{inc:AN02:tail:v1:sources}
\subsection{Ancestry must not change when a label changes chart}
One possible disjoint initial partition, up to null coordinate axes, is
Q1, Q2, Q3, Q4. Q1 labels are candidates for strong-tail recruitment, Q2
labels for bulk weak arrival, and Q3 labels for the reservoir. These names
specify \emph{sources}, not proved future specialization. Q4 and any
uncaptured portions remain in the accounting.
For a fixed captured label set $C$ and an ancestral source $B_j$, define
\[
 N_j(\tau)=\int_{C\cap B_j}m_\tau\dd\lambda_0.
\]
In the original flow,
\begin{equation}\label{inc:AN02:tail:v1:exactsource}
 N_j'=\int_{C\cap B_j}2b^\top R_f(\theta)a\,m\dd\lambda_0.
\end{equation}
Finite-horizon characteristic bounds justify differentiation. No moving
membership is hidden in this formula; whether $C$ stays in the weak chart
is an additional requirement.

\begin{proposition}[Closed-source laws and recruitment defects]
\label{inc:AN02:tail:v1:sourceclock}
Suppose absolutely continuous nonnegative source masses $N_j$ have
$N=\sum_jN_j\in(0,1)$ and obey
\begin{equation}\label{inc:AN02:tail:v1:fluxlaw}
 N_j'=\lambda(1-N)N_j+j_j,\qquad j=\sum_jj_j.
\end{equation}
Then
\begin{equation}\label{inc:AN02:tail:v1:fluxidentities}
 N'=\lambda N(1-N)+j,\qquad
 \left(\frac{N_j}N\right)'=\frac{j_j-(N_j/N)j}{N},\qquad
 \left(\log\frac N{1-N}\right)'=\lambda+\frac{j}{N(1-N)}.
\end{equation}
For a closed source decomposition of the leading weak family, $j_j=0$:
source proportions are constant and total mass is exactly logistic.
For a moving weak mask, $j_j$ must include its entry/exit flux as well as
any finite-noise reaction defect. It cannot be discarded by changing a
label's name from strong to weak.
\end{proposition}
\begin{proof}
Sum \eqref{inc:AN02:tail:v1:fluxlaw}, differentiate $N_j/N$, and use
$(\log[N/(1-N)])'=N'/[N(1-N)]$. In the closed leading family every
source has reaction rate $\lambda(1-N)$ by the defining transport equation.
For the original finite-noise flow on a fixed set, the defect is exactly
\[
 j_j=\int_{C\cap B_j}
 [2b^\top R_fa-\lambda(1-N)]m\dd\lambda_0.
\]
The proposition is an identity for specified absolutely continuous masses;
it does not assume or prove that a particular moving mask has an
absolutely continuous flux. Singular boundary crossings require the
corresponding measure-valued bookkeeping instead.
\end{proof}

\subsection{The bounded interval is the clock-to-witness hull}
Keep AN06 v2's fixed shifted witness interval $J$ and define
\[
 \mathcal K_J=\operatorname{conv}(\{0\}\cup J),\qquad
 L_J=\max_{s\in\mathcal K_J}|s|.
\]
Center at the \emph{population's own} weak half-mass time $\tau_h$, so
$s=\tau-\tau_h$ and $N(0)=1/2$. For the closed leading family,
$N(s)=n_\rho(s)=(1+e^{-\lambda s})^{-1}$ exactly. The AN06 distance must
therefore use its phase-aligned form, omitting $|N-n_\rho|$.

If the original reduced family instead has mass defect $r_N$,
$N'=\lambda N(1-N)+r_N$, then the reviewed logit identity gives
\begin{equation}\label{inc:AN02:tail:v1:windowclock}
 |N(s)-n_\rho(s)|\le\frac14
 \left|\int_0^s\frac{r_N(t)}{N(t)(1-N(t))}\dd t\right|.
\end{equation}
In particular, if $\nu\le N\le1-\nu$ on $\mathcal K_J$, $0<\nu<1/2$,
\begin{equation}\label{inc:AN02:tail:v1:windowclockbound}
 \sup_{s\in J}|N(s)-n_\rho(s)|
 \le\frac{L_J}{4\nu(1-\nu)}\|r_N\|_{L^\infty(\mathcal K_J)}.
\end{equation}
The uniform separation from one is part of the bounded-window guard;
merely writing $N<1$ is not a quantitative denominator bound. Since the
coherent logistic orbit on a compact parameter/time set has such a guard,
a sufficiently small defect closes it by the usual first-exit argument.
The hull, rather than $J$ alone, is necessary when the half-mass time is
outside $J$.

The intended AN3/AN7 interface is thus as follows. Long-delay entry and
passage must control the strong mean, normalized shapes and tails, weak
angular/source law, and outside population in the population's half-mass
clock. Long-sojourn relative reaction errors, including
$q^2\log(1/S_0)$, enter those state/profile estimates and the control of the
half-mass event; they are not charged a second time as an early-centered
phase mismatch on $J$. The local mass-law defect is integrated only over
$\mathcal K_J$. Non-negligible late recruitment contributes to that defect
and to source proportions through \eqref{inc:AN02:tail:v1:fluxidentities}.

\paragraph{AN06 probe caps.}
The endpoints restricting the local probe rectangle to $[1.74,1.75]$ in
AN06 v2 are sufficient containment choices, not a necessary dynamical
restriction. They can be removed in a later authorized consolidation,
retaining the central strict margins and the actual robustness bounds.
No version is overwritten for this cosmetic change.

\section{What the coupled initialized lemma must still supply}
\label{inc:AN02:tail:v1:remaining}
This increment provides a genuine original-flow ancestral mass lower
bound and two different analytical tools for tail motion. It does not
establish the strong mean/shape entry, the intermediate angular matching,
weak capture, or finite-noise resident passage.

The next useful output is not just a variance at a putative saturation
state. It must specify an actual same-field core characteristic or a
controlled core reference; the strong mean; an ancestral-to-spatial tail
pushforward with upper as well as lower control; and a source-resolved
weak-entry law. The tail estimate must then be propagated over the full
resident interval, for example through
\eqref{inc:AN02:tail:v1:exitmass}, with every exiting label retained in a
subsequent chart or outside-population budget.

In particular the following implications remain unproved:
\[
 \begin{aligned}
 \text{original ancestral weight}
 &\ \Longrightarrow\ \text{spatial profile at learned strong entry}\\
 &\ \Longrightarrow\ \text{controlled recruitment/capture}.
 \end{aligned}
\]
The source labels Q1/Q2/Q3 cannot yet be assigned either a proved relative
contribution or a dominant weak-seed exponent. The corrected $5/6$ outer
proxy makes the proposed dismissal of recruitment unavailable. Conversely,
it does not prove that recruitment dominates the actual initialized seed.

The local-core $\alpha_{\rm sh}$ rate and the global-tail $\lambda/2$
rate must not be interchanged in the AN3 power balance. No positivity claim
for a final combined exponent is made without the initialized profile,
seed, and finite-noise estimates. A theorem in the deep joint regime
$\log(1/S_0)/\log(1/q)\to\infty$ and
$q^2\log(1/S_0)\to0$ remains a legitimate target, but neither condition
alone supplies these missing estimates.

\begin{thebibliography}{9}
\raggedright
\bibitem{F} \path{RELU_POPULATION_FOUNDATIONS_v2.tex}.
Frozen supplied foundations; P4 and P5 for Gaussian/characteristic formulas.
\bibitem{P} \path{THEOREM_A_ANALYTICAL_PROGRESS.tex}, version 1.0.
Labels \texttt{pa:earlybounds}, \texttt{pa:targetcomparison},
\texttt{pa:overlap}, \texttt{pa:leading}, \texttt{pa:masslaws},
\texttt{pa:resident}, \texttt{pa:order}, and \texttt{pa:multipliers}.
\bibitem{A6} \path{AN06_learned_orbit_witnesses_v2.tex}.
Reviewed canonical-offset learned witnesses and phase-aligned population tolerance.
\bibitem{H} \path{THEORY_HANDOFF_ANALYTICAL_A_B.md}.
Analytical-only proof contract and separate-review-file workflow.
\bibitem{R} User's review of AN06 v2 in the current conversation.
Motivation for the weighted ancestry profile, resident tail propagation,
source splitting, and half-mass-clock clarification. The exponent and the
far-field separatrix interpretation are corrected explicitly above.
\end{thebibliography}
\end{document}
